% ============================================================
%  PANEL DE CONTROL — configuracion.tex
%  ------------------------------------------------------------
%  TODO lo que cambia de una versión a otra vive AQUÍ. Ni
%  main.tex ni los archivos de secciones/ llevan datos ni
%  interruptores de formato.
%
%  Uso:  main.tex empieza con  % ============================================================
%  PANEL DE CONTROL — configuracion.tex
%  ------------------------------------------------------------
%  TODO lo que cambia de una versión a otra vive AQUÍ. Ni
%  main.tex ni los archivos de secciones/ llevan datos ni
%  interruptores de formato.
%
%  Uso:  main.tex empieza con  % ============================================================
%  PANEL DE CONTROL — configuracion.tex
%  ------------------------------------------------------------
%  TODO lo que cambia de una versión a otra vive AQUÍ. Ni
%  main.tex ni los archivos de secciones/ llevan datos ni
%  interruptores de formato.
%
%  Uso:  main.tex empieza con  % ============================================================
%  PANEL DE CONTROL — configuracion.tex
%  ------------------------------------------------------------
%  TODO lo que cambia de una versión a otra vive AQUÍ. Ni
%  main.tex ni los archivos de secciones/ llevan datos ni
%  interruptores de formato.
%
%  Uso:  main.tex empieza con  \input{configuracion.tex}
%        (antes del preámbulo, para que este pueda reaccionar a
%        los interruptores) y después entra el contenido.
% ============================================================

% ============================================================
%  1. DATOS DEL PROYECTO
% ============================================================

% ---------- Título de la monografía ------------------------
% El título se imprime en MAYÚSCULAS, Arial 20 negrita, centrado
% y con los saltos de línea ESCRITOS A MANO: la guía no permite
% reducir el tamaño de la fuente, así que la pirámide invertida
% se consigue cortando el título en renglones explícitos.
% El bloque útil mide 16,09 cm; a Arial 20 negrita el renglón más
% ancho que puede empezar por "IDENTIFICACIÓN" mide 11,8 cm
% (la palabra siguiente ya no cabe), de modo que la pirámide
% necesariamente tiene su base en los renglones 2 y 3 y se
% estrecha hacia abajo, que es la silueta del modelo.
% Anchuras de los renglones, de arriba abajo:
%   11,8 · 15,8 · 15,3 · 12,0 · 12,6 · 3,8 cm
% Para cambiar el corte basta reescribir estas líneas: si un
% renglón no cabe en 16,09 cm se partirá en dos (y la pirámide se
% pierde), así que conviene medir el resultado con
% "scripts/verificar" tras cada ajuste.

% Interruptor: modo borrador con figuras pendientes visibles (true)
% o modo entrega sin recuadros (false). Ver preambulo.tex sección 12.
\newif\ifborrador
\borradortrue  % default: muestra recuadros de "Figura pendiente"
\newcommand{\tituloLineaA}{Estimación del número esperado de}
\newcommand{\tituloLineaB}{consultas de ruta de Trufi App por}
\newcommand{\tituloLineaC}{celda H3 en el eje metropolitano}
\newcommand{\tituloLineaD}{de Cochabamba mediante técnicas}
\newcommand{\tituloLineaE}{geoespaciales de Ciencia de Datos}
\newcommand{\tituloMonografia}{%
  \tituloLineaA\\
  \tituloLineaB\\
  \tituloLineaC\\
  \tituloLineaD\\
  \tituloLineaE}
% Versión en una sola línea, solo para los metadatos del PDF
% (\hypersetup no admite saltos de línea).
\newcommand{\tituloMonografiaPlano}{Estimación del número esperado de consultas de ruta
  de Trufi App por celda H3 en el eje metropolitano de Cochabamba mediante técnicas
  geoespaciales de Ciencia de Datos}
% Título corto: el del encabezado corrido del cuerpo (Fase 5).
\newcommand{\tituloCorto}{Demanda esperada de consultas de Trufi App por celda H3}

% ---------- Postulante, programa, versión y año --------------
\newcommand{\nombreDiplomado}{Ciencia de Datos}
\newcommand{\nombreVersion}{CUARTA}
\newcommand{\nombrePostulante}{José Roberto Vargas Orellana}
\newcommand{\anioMonografia}{2026}
\newcommand{\ciudadMonografia}{Cochabamba -- Bolivia}

% ---------- Docente y módulo (línea oculta por defecto) -----
% La guía los menciona en §2 pero no aparecen en el modelo de
% la carátula, así que quedan ocultos. Para mostrarlos basta
% poner \mostrarDocentetrue más abajo.
\newcommand{\nombreDocente}{Nombre del docente}
\newcommand{\nombreModulo}{Nombre del módulo}

% ============================================================
%  2. INTERRUPTORES
% ------------------------------------------------------------
%  El valor por defecto de cada línea es el de la guía; la
%  alternativa se consigue cambiando solo el "true"/"false".
% ============================================================

% ---------- Numeración de páginas ---------------------------
%   false (guía)  = romanos en los preliminares + arábigos desde 1
%   true  (modelo)= arábiga continua desde la carátula
\newif\ifnumeracionContinua
\numeracionContinuafalse

% ---------- Posición del folio ------------------------------
%   false = abajo y centrado (lo exige la guía)
%   true  = abajo a la derecha
\newif\iffolioDerecha
\folioDerechafalse

% ---------- Línea docente/módulo en la carátula --------------
%   false = oculta · true = visible
\newif\ifmostrarDocente
\mostrarDocentefalse

% ---------- Modo del documento ------------------------------
%   false = final      (falla la compilación si quedan marcadores)
%   true  = borrador   (resalta los marcadores de trabajo)
\newif\ifmodoBorrador
\modoBorradorfalse

% ---------- Color ------------------------------------------
%   true  = digital, con el azul marino de acento
%   false = impresión, todo en negro y gris
\newif\ifcolorDigital
\colorDigitaltrue

% ---------- Preliminares opcionales ------------------------
%   true  = mostrar dedicatoria, agradecimientos, glosario e
%           índices de tablas y figuras si existen
%   false = ocultarlos
\newif\ifmostrarOpcionales
\mostrarOpcionalestrue

% ---------- Compilar un solo capítulo ----------------------
%   false = documento completo
%   true  = solo el capítulo indicado en \capituloUnico
%           (útil para escribir una sección sin recompilar todo)
\newif\ifcompilarUnCapitulo
\compilarUnCapitulofalse
\newcommand{\capituloUnico}{08_desarrollo}

% ---------- Estilo de bibliografía -------------------------
%   false = APA: alfabética, sin numerar (predeterminado)
%   true  = numerada, formato de la guía §3.6
\newif\ifbibNumerada
\bibNumeradafalse

%        (antes del preámbulo, para que este pueda reaccionar a
%        los interruptores) y después entra el contenido.
% ============================================================

% ============================================================
%  1. DATOS DEL PROYECTO
% ============================================================

% ---------- Título de la monografía ------------------------
% El título se imprime en MAYÚSCULAS, Arial 20 negrita, centrado
% y con los saltos de línea ESCRITOS A MANO: la guía no permite
% reducir el tamaño de la fuente, así que la pirámide invertida
% se consigue cortando el título en renglones explícitos.
% El bloque útil mide 16,09 cm; a Arial 20 negrita el renglón más
% ancho que puede empezar por "IDENTIFICACIÓN" mide 11,8 cm
% (la palabra siguiente ya no cabe), de modo que la pirámide
% necesariamente tiene su base en los renglones 2 y 3 y se
% estrecha hacia abajo, que es la silueta del modelo.
% Anchuras de los renglones, de arriba abajo:
%   11,8 · 15,8 · 15,3 · 12,0 · 12,6 · 3,8 cm
% Para cambiar el corte basta reescribir estas líneas: si un
% renglón no cabe en 16,09 cm se partirá en dos (y la pirámide se
% pierde), así que conviene medir el resultado con
% "scripts/verificar" tras cada ajuste.

% Interruptor: modo borrador con figuras pendientes visibles (true)
% o modo entrega sin recuadros (false). Ver preambulo.tex sección 12.
\newif\ifborrador
\borradortrue  % default: muestra recuadros de "Figura pendiente"
\newcommand{\tituloLineaA}{Estimación del número esperado de}
\newcommand{\tituloLineaB}{consultas de ruta de Trufi App por}
\newcommand{\tituloLineaC}{celda H3 en el eje metropolitano}
\newcommand{\tituloLineaD}{de Cochabamba mediante técnicas}
\newcommand{\tituloLineaE}{geoespaciales de Ciencia de Datos}
\newcommand{\tituloMonografia}{%
  \tituloLineaA\\
  \tituloLineaB\\
  \tituloLineaC\\
  \tituloLineaD\\
  \tituloLineaE}
% Versión en una sola línea, solo para los metadatos del PDF
% (\hypersetup no admite saltos de línea).
\newcommand{\tituloMonografiaPlano}{Estimación del número esperado de consultas de ruta
  de Trufi App por celda H3 en el eje metropolitano de Cochabamba mediante técnicas
  geoespaciales de Ciencia de Datos}
% Título corto: el del encabezado corrido del cuerpo (Fase 5).
\newcommand{\tituloCorto}{Demanda esperada de consultas de Trufi App por celda H3}

% ---------- Postulante, programa, versión y año --------------
\newcommand{\nombreDiplomado}{Ciencia de Datos}
\newcommand{\nombreVersion}{CUARTA}
\newcommand{\nombrePostulante}{José Roberto Vargas Orellana}
\newcommand{\anioMonografia}{2026}
\newcommand{\ciudadMonografia}{Cochabamba -- Bolivia}

% ---------- Docente y módulo (línea oculta por defecto) -----
% La guía los menciona en §2 pero no aparecen en el modelo de
% la carátula, así que quedan ocultos. Para mostrarlos basta
% poner \mostrarDocentetrue más abajo.
\newcommand{\nombreDocente}{Nombre del docente}
\newcommand{\nombreModulo}{Nombre del módulo}

% ============================================================
%  2. INTERRUPTORES
% ------------------------------------------------------------
%  El valor por defecto de cada línea es el de la guía; la
%  alternativa se consigue cambiando solo el "true"/"false".
% ============================================================

% ---------- Numeración de páginas ---------------------------
%   false (guía)  = romanos en los preliminares + arábigos desde 1
%   true  (modelo)= arábiga continua desde la carátula
\newif\ifnumeracionContinua
\numeracionContinuafalse

% ---------- Posición del folio ------------------------------
%   false = abajo y centrado (lo exige la guía)
%   true  = abajo a la derecha
\newif\iffolioDerecha
\folioDerechafalse

% ---------- Línea docente/módulo en la carátula --------------
%   false = oculta · true = visible
\newif\ifmostrarDocente
\mostrarDocentefalse

% ---------- Modo del documento ------------------------------
%   false = final      (falla la compilación si quedan marcadores)
%   true  = borrador   (resalta los marcadores de trabajo)
\newif\ifmodoBorrador
\modoBorradorfalse

% ---------- Color ------------------------------------------
%   true  = digital, con el azul marino de acento
%   false = impresión, todo en negro y gris
\newif\ifcolorDigital
\colorDigitaltrue

% ---------- Preliminares opcionales ------------------------
%   true  = mostrar dedicatoria, agradecimientos, glosario e
%           índices de tablas y figuras si existen
%   false = ocultarlos
\newif\ifmostrarOpcionales
\mostrarOpcionalestrue

% ---------- Compilar un solo capítulo ----------------------
%   false = documento completo
%   true  = solo el capítulo indicado en \capituloUnico
%           (útil para escribir una sección sin recompilar todo)
\newif\ifcompilarUnCapitulo
\compilarUnCapitulofalse
\newcommand{\capituloUnico}{08_desarrollo}

% ---------- Estilo de bibliografía -------------------------
%   false = APA: alfabética, sin numerar (predeterminado)
%   true  = numerada, formato de la guía §3.6
\newif\ifbibNumerada
\bibNumeradafalse

%        (antes del preámbulo, para que este pueda reaccionar a
%        los interruptores) y después entra el contenido.
% ============================================================

% ============================================================
%  1. DATOS DEL PROYECTO
% ============================================================

% ---------- Título de la monografía ------------------------
% El título se imprime en MAYÚSCULAS, Arial 20 negrita, centrado
% y con los saltos de línea ESCRITOS A MANO: la guía no permite
% reducir el tamaño de la fuente, así que la pirámide invertida
% se consigue cortando el título en renglones explícitos.
% El bloque útil mide 16,09 cm; a Arial 20 negrita el renglón más
% ancho que puede empezar por "IDENTIFICACIÓN" mide 11,8 cm
% (la palabra siguiente ya no cabe), de modo que la pirámide
% necesariamente tiene su base en los renglones 2 y 3 y se
% estrecha hacia abajo, que es la silueta del modelo.
% Anchuras de los renglones, de arriba abajo:
%   11,8 · 15,8 · 15,3 · 12,0 · 12,6 · 3,8 cm
% Para cambiar el corte basta reescribir estas líneas: si un
% renglón no cabe en 16,09 cm se partirá en dos (y la pirámide se
% pierde), así que conviene medir el resultado con
% "scripts/verificar" tras cada ajuste.

% Interruptor: modo borrador con figuras pendientes visibles (true)
% o modo entrega sin recuadros (false). Ver preambulo.tex sección 12.
\newif\ifborrador
\borradortrue  % default: muestra recuadros de "Figura pendiente"
\newcommand{\tituloLineaA}{Estimación del número esperado de}
\newcommand{\tituloLineaB}{consultas de ruta de Trufi App por}
\newcommand{\tituloLineaC}{celda H3 en el eje metropolitano}
\newcommand{\tituloLineaD}{de Cochabamba mediante técnicas}
\newcommand{\tituloLineaE}{geoespaciales de Ciencia de Datos}
\newcommand{\tituloMonografia}{%
  \tituloLineaA\\
  \tituloLineaB\\
  \tituloLineaC\\
  \tituloLineaD\\
  \tituloLineaE}
% Versión en una sola línea, solo para los metadatos del PDF
% (\hypersetup no admite saltos de línea).
\newcommand{\tituloMonografiaPlano}{Estimación del número esperado de consultas de ruta
  de Trufi App por celda H3 en el eje metropolitano de Cochabamba mediante técnicas
  geoespaciales de Ciencia de Datos}
% Título corto: el del encabezado corrido del cuerpo (Fase 5).
\newcommand{\tituloCorto}{Demanda esperada de consultas de Trufi App por celda H3}

% ---------- Postulante, programa, versión y año --------------
\newcommand{\nombreDiplomado}{Ciencia de Datos}
\newcommand{\nombreVersion}{CUARTA}
\newcommand{\nombrePostulante}{José Roberto Vargas Orellana}
\newcommand{\anioMonografia}{2026}
\newcommand{\ciudadMonografia}{Cochabamba -- Bolivia}

% ---------- Docente y módulo (línea oculta por defecto) -----
% La guía los menciona en §2 pero no aparecen en el modelo de
% la carátula, así que quedan ocultos. Para mostrarlos basta
% poner \mostrarDocentetrue más abajo.
\newcommand{\nombreDocente}{Nombre del docente}
\newcommand{\nombreModulo}{Nombre del módulo}

% ============================================================
%  2. INTERRUPTORES
% ------------------------------------------------------------
%  El valor por defecto de cada línea es el de la guía; la
%  alternativa se consigue cambiando solo el "true"/"false".
% ============================================================

% ---------- Numeración de páginas ---------------------------
%   false (guía)  = romanos en los preliminares + arábigos desde 1
%   true  (modelo)= arábiga continua desde la carátula
\newif\ifnumeracionContinua
\numeracionContinuafalse

% ---------- Posición del folio ------------------------------
%   false = abajo y centrado (lo exige la guía)
%   true  = abajo a la derecha
\newif\iffolioDerecha
\folioDerechafalse

% ---------- Línea docente/módulo en la carátula --------------
%   false = oculta · true = visible
\newif\ifmostrarDocente
\mostrarDocentefalse

% ---------- Modo del documento ------------------------------
%   false = final      (falla la compilación si quedan marcadores)
%   true  = borrador   (resalta los marcadores de trabajo)
\newif\ifmodoBorrador
\modoBorradorfalse

% ---------- Color ------------------------------------------
%   true  = digital, con el azul marino de acento
%   false = impresión, todo en negro y gris
\newif\ifcolorDigital
\colorDigitaltrue

% ---------- Preliminares opcionales ------------------------
%   true  = mostrar dedicatoria, agradecimientos, glosario e
%           índices de tablas y figuras si existen
%   false = ocultarlos
\newif\ifmostrarOpcionales
\mostrarOpcionalestrue

% ---------- Compilar un solo capítulo ----------------------
%   false = documento completo
%   true  = solo el capítulo indicado en \capituloUnico
%           (útil para escribir una sección sin recompilar todo)
\newif\ifcompilarUnCapitulo
\compilarUnCapitulofalse
\newcommand{\capituloUnico}{08_desarrollo}

% ---------- Estilo de bibliografía -------------------------
%   false = APA: alfabética, sin numerar (predeterminado)
%   true  = numerada, formato de la guía §3.6
\newif\ifbibNumerada
\bibNumeradafalse

%        (antes del preámbulo, para que este pueda reaccionar a
%        los interruptores) y después entra el contenido.
% ============================================================

% ============================================================
%  1. DATOS DEL PROYECTO
% ============================================================

% ---------- Título de la monografía ------------------------
% El título se imprime en MAYÚSCULAS, Arial 20 negrita, centrado
% y con los saltos de línea ESCRITOS A MANO: la guía no permite
% reducir el tamaño de la fuente, así que la pirámide invertida
% se consigue cortando el título en renglones explícitos.
% El bloque útil mide 16,09 cm; a Arial 20 negrita el renglón más
% ancho que puede empezar por "IDENTIFICACIÓN" mide 11,8 cm
% (la palabra siguiente ya no cabe), de modo que la pirámide
% necesariamente tiene su base en los renglones 2 y 3 y se
% estrecha hacia abajo, que es la silueta del modelo.
% Anchuras de los renglones, de arriba abajo:
%   11,8 · 15,8 · 15,3 · 12,0 · 12,6 · 3,8 cm
% Para cambiar el corte basta reescribir estas líneas: si un
% renglón no cabe en 16,09 cm se partirá en dos (y la pirámide se
% pierde), así que conviene medir el resultado con
% "scripts/verificar" tras cada ajuste.

% Interruptor: modo borrador con figuras pendientes visibles (true)
% o modo entrega sin recuadros (false). Ver preambulo.tex sección 12.
\newif\ifborrador
\borradortrue  % default: muestra recuadros de "Figura pendiente"
\newcommand{\tituloLineaA}{Estimación del número esperado de}
\newcommand{\tituloLineaB}{consultas de ruta de Trufi App por}
\newcommand{\tituloLineaC}{celda H3 en el eje metropolitano}
\newcommand{\tituloLineaD}{de Cochabamba mediante técnicas}
\newcommand{\tituloLineaE}{geoespaciales de Ciencia de Datos}
\newcommand{\tituloMonografia}{%
  \tituloLineaA\\
  \tituloLineaB\\
  \tituloLineaC\\
  \tituloLineaD\\
  \tituloLineaE}
% Versión en una sola línea, solo para los metadatos del PDF
% (\hypersetup no admite saltos de línea).
\newcommand{\tituloMonografiaPlano}{Estimación del número esperado de consultas de ruta
  de Trufi App por celda H3 en el eje metropolitano de Cochabamba mediante técnicas
  geoespaciales de Ciencia de Datos}
% Título corto: el del encabezado corrido del cuerpo (Fase 5).
\newcommand{\tituloCorto}{Demanda esperada de consultas de Trufi App por celda H3}

% ---------- Postulante, programa, versión y año --------------
\newcommand{\nombreDiplomado}{Ciencia de Datos}
\newcommand{\nombreVersion}{CUARTA}
\newcommand{\nombrePostulante}{José Roberto Vargas Orellana}
\newcommand{\anioMonografia}{2026}
\newcommand{\ciudadMonografia}{Cochabamba -- Bolivia}

% ---------- Docente y módulo (línea oculta por defecto) -----
% La guía los menciona en §2 pero no aparecen en el modelo de
% la carátula, así que quedan ocultos. Para mostrarlos basta
% poner \mostrarDocentetrue más abajo.
\newcommand{\nombreDocente}{Nombre del docente}
\newcommand{\nombreModulo}{Nombre del módulo}

% ============================================================
%  2. INTERRUPTORES
% ------------------------------------------------------------
%  El valor por defecto de cada línea es el de la guía; la
%  alternativa se consigue cambiando solo el "true"/"false".
% ============================================================

% ---------- Numeración de páginas ---------------------------
%   false (guía)  = romanos en los preliminares + arábigos desde 1
%   true  (modelo)= arábiga continua desde la carátula
\newif\ifnumeracionContinua
\numeracionContinuafalse

% ---------- Posición del folio ------------------------------
%   false = abajo y centrado (lo exige la guía)
%   true  = abajo a la derecha
\newif\iffolioDerecha
\folioDerechafalse

% ---------- Línea docente/módulo en la carátula --------------
%   false = oculta · true = visible
\newif\ifmostrarDocente
\mostrarDocentefalse

% ---------- Modo del documento ------------------------------
%   false = final      (falla la compilación si quedan marcadores)
%   true  = borrador   (resalta los marcadores de trabajo)
\newif\ifmodoBorrador
\modoBorradorfalse

% ---------- Color ------------------------------------------
%   true  = digital, con el azul marino de acento
%   false = impresión, todo en negro y gris
\newif\ifcolorDigital
\colorDigitaltrue

% ---------- Preliminares opcionales ------------------------
%   true  = mostrar dedicatoria, agradecimientos, glosario e
%           índices de tablas y figuras si existen
%   false = ocultarlos
\newif\ifmostrarOpcionales
\mostrarOpcionalestrue

% ---------- Compilar un solo capítulo ----------------------
%   false = documento completo
%   true  = solo el capítulo indicado en \capituloUnico
%           (útil para escribir una sección sin recompilar todo)
\newif\ifcompilarUnCapitulo
\compilarUnCapitulofalse
\newcommand{\capituloUnico}{08_desarrollo}

% ---------- Estilo de bibliografía -------------------------
%   false = APA: alfabética, sin numerar (predeterminado)
%   true  = numerada, formato de la guía §3.6
\newif\ifbibNumerada
\bibNumeradafalse
